\chapter{Questionário com frequentadores}
\label{ap:questionario}

Este apêndice apresenta o questionário a ser aplicado a frequentadores de bares e restaurantes de Porto Alegre. O instrumento será hospedado no Google Forms e distribuído por meio de redes sociais, grupos de mensageria temáticos, listas institucionais e, opcionalmente, por QR Code disponibilizado nos estabelecimentos participantes do estudo de caso. O tempo de resposta estimado é de três a cinco minutos. A meta de coleta é de no mínimo 80 respostas, com 150 ou mais como meta desejável. O questionário inicia com uma versão sumarizada do Termo de Consentimento Livre e Esclarecido (Apêndice~\ref{ap:termos}), exigindo concordância antes do prosseguimento.

\section*{Seção 1 --- Perfil do respondente}

\begin{enumerate}
	\item Faixa etária: menor de 18; de 18 a 24; de 25 a 34; de 35 a 44; de 45 a 54; 55 ou mais.
	\item Você mora em Porto Alegre ou na região metropolitana? Sim; Não.
	\item Com que frequência você frequenta bares ou restaurantes? Várias vezes por semana; uma vez por semana; quinzenalmente; mensalmente; raramente.
	\item Que tipo de estabelecimento você frequenta com maior frequência? Bar ou boteco; restaurante casual; restaurante sofisticado; \emph{fast-casual}; outro.
	\item Você possui alguma deficiência ou dificuldade que afete o uso de cardápios? Visual; motora; de leitura; nenhuma; prefiro não responder.
\end{enumerate}

\section*{Seção 2 --- Experiência com cardápios digitais}

\begin{enumerate}\setcounter{enumi}{5}
	\item Nos últimos 12 meses, com que frequência você se deparou com cardápios digitais (QR Code ou \emph{web app}) ao frequentar bares e restaurantes? Sempre; frequentemente; às vezes; raramente; nunca. \emph{(Em caso de ``nunca'', avançar diretamente para a Seção~4.)}
	\item Em geral, qual é a sua preferência? Cardápio físico; cardápio digital; indiferente.
	\item Indique seu nível de concordância com as afirmações a seguir, em escala de 1 (discordo totalmente) a 5 (concordo totalmente):
	\begin{alineas}
		\item ``cardápio digital é mais prático que o físico'';
		\item ``tenho dificuldade de ler ou navegar no cardápio digital'';
		\item ``sinto falta de interação com o garçom quando uso o cardápio digital'';
		\item ``confio mais nos preços e na disponibilidade de itens informados pelo cardápio digital'';
		\item ``me preocupo com a proteção dos meus dados ao usar QR Code em restaurantes'';
		\item ``o cardápio digital me leva a pedir mais ou a gastar mais'';
		\item ``o cardápio digital atrapalha a experiência social com as pessoas que estão comigo''.
	\end{alineas}
	\item Você já desistiu de pedir ou ficou frustrado em alguma situação envolvendo cardápio digital? Sim; Não. \emph{(Em caso afirmativo, descreva brevemente.)}
\end{enumerate}

\section*{Seção 3 --- Aspectos técnicos}

\begin{enumerate}\setcounter{enumi}{9}
	\item Você já enfrentou algum dos problemas a seguir? (marcação múltipla) QR Code não foi reconhecido pela câmera; página não carregou; \emph{layout} ruim ou confuso; preço ou informação desatualizados; nenhum desses.
	\item Em que situações o cardápio digital funciona melhor para você? \emph{(Resposta aberta, curta.)}
	\item Em que situações o cardápio digital funciona pior para você? \emph{(Resposta aberta, curta.)}
\end{enumerate}

\section*{Seção 4 --- Visão geral}

\begin{enumerate}\setcounter{enumi}{12}
	\item Você acredita que os cardápios digitais vieram para ficar? Sim; Não; Não sei.
	\item Se você fosse dono de um bar ou restaurante, adotaria cardápio digital? Sim; Não; Talvez. \emph{(Justificativa opcional.)}
	\item Há algo mais que você gostaria de comentar sobre o tema? \emph{(Resposta aberta, opcional.)}
\end{enumerate}
